% English translation of capitulos/04a3-traducao-refatoracao.tex, made by Claude Code (Claude Fable 5.2) on 24 November 2026.
% Pending review by the author. When the Portuguese file changes, update this one (scripts/verificar_traducao.py).
\subsubsection{Translation}

The partial synthesis of translation covers 25 code translation studies, and the 2 included later are at the end of this
part; none of them also records the refactoring task. The records that support this subsection were extracted by AI and checked by the author only in a sample; the claims were checked against the original. Counting each study once per type of counterfactual, with overlap between the types, 17 of the 25 studies evaluate at least one package, 12 isolate some ingredient, 6 have a cumulative design and 3 (S1010, S1220 and S1836) have no counterfactual. The most frequent sets are CodeNet (6 of the 25 studies), TransCoder/GeeksforGeeks (5) and AVATAR (5). By the mechanism added to the prompt, the studies form six groupings:
\begin{itemize}
  \item repair with feedback from compilation, execution or tests (11 studies);
  \item examples, in quantity, selection or retrieval (6);
  \item explanation, reasoning and intermediate representation (2);
  \item instructions, constraints and context oriented to requirements (6);
  \item automatic prompt optimization (2, of which 1 with a translation comparison);
  \item hybrid systems and studies without a counterfactual (4).
\end{itemize}

\textbf{Repair with execution feedback.} In 9 of the 11 studies, the focal technique is evaluated as a package; only S329 and S1727 isolate ingredients by ablation.
% AFIRMACAO: A25
Among the packages, computational accuracy (CA) went from 54\% to 80\% in S732 (GPT-4o, from C to Rust), the syntax pass rate went from 51.9\% to 75.0\% in S1132 and the pass rate went from 80\% to 97\% in S2118, against an agent without compiler analysis, although in the latter the median speedup of the package was 0.76$\times$, against 0.88$\times$ for the condition without analysis (S732-C1; S1132-C1; S2118-C1; S2118).
% AFIRMACAO: A26
In the opposite direction, in S2051 with CodeLlama-70b, the pass rate went from 29.71\% (initial pass@3) to 16.57\% after the first repair iteration, measurements that are not equivalent, and the selection of the best version among the initial and the repaired ones reached 39.97\% (S2051-C1; S2051). In S919, Pass went from 89.5\% to 89.2\% (S919-C3), and in S1163 compilation stayed at the ceiling in both conditions in the nanoXOR cell without additional context (S1163-C1).

% AFIRMACAO: A27
In S329, with LLaMA-7B, the package raised CA from 22.74\% to 28.45\%; in the ablations, the test cases in the prompt took CA from 22.74\% to 27.43\%, and execution-guided repair, added to them, from 27.43\% to 28.45\% (S329-C1; S329-C2; S329-C3).
One possible reading is that, in this model and set, most of the difference of the package accompanies the addition of the tests, and not the repair. The studies also differ in the number of rounds (up to three in S1727, up to five in S732 and up to ten in S1132), and none measures cost, although the number of calls to the model varies between them.

\textbf{Examples.} The six studies compare \textit{few-shot}, \textit{many-shot} and selection of examples by similarity or structure.
% AFIRMACAO: A28
With examples, CA went from 61.28\% to 64.28\% in S674 (Qwen2-7B), with structural against random selection and five examples in both conditions, BLEU went from 12.14 to 12.58 in S1032 with five examples and the syntax pass rate went from 51.9\% to 60.0\% in S1132 with simple retrieval of examples (S674-C1; S1032-C1; S1132-C3).
% AFIRMACAO: A29
In the opposite direction, the prompt with one example reduced the compilation rate from 0.49 to 0.43 in S964 with LLaMA-3.3-70B (S964-C1), and in S409, on another metric (CodeBLEU), StarCoder 15.5B did not vary (delta of 0.000) with one, two or three examples (S409).
These results suggest, without isolating the factor, dependence on the model and on the selection criterion. In addition, S409 and S1032 use surface metrics (BLEU and CodeBLEU), which do not measure functional correctness.

\textbf{Explanation and reasoning.}
% AFIRMACAO: A30
In S193, with GPT-3.5 from Python to JavaScript, the three variants of explanation before the translation had pass@1 of 84.8\%, 84.3\% and 83.2\%, against 85.5\% for direct translation, although the introduction of the study claims consistent improvement (S193-C1; S193-C2; S193-C3; S193).
% AFIRMACAO: A31
In S836 (Open GPT4 8X7B, CodeNet), the \textit{one-shot} prompts with reasoning raised the success rate from 28.6\% to values between 32.2\% and 42.4\%, whereas the four two-step \textit{zero-shot} approaches (code, then natural language or AST, then target language) reduced it to 10.7\%, 2.6\%, 9.0\% and 11.2\% (S836-C1; S836-C4; S836-C7; S836-C8; S836-C2; S836-C3; S836-C5; S836-C6).
Since the conditions with reasoning also add an example, this gain is a package gain and cannot be attributed to the reasoning.

\textbf{Instructions and requirements.}
% AFIRMACAO: A32
In S304, with CodeLlama 7B and the same regular expression in both conditions, the format constraint raised the output extraction rate (MSR) from 6.6\% to 94.2\%, but in Mixtral the relation was reversed (MSR of 60.4\% against 82.20\%; CA of 0\% against 35.60\%) (S304-C1; S304).
MSR measures the extraction of the answer, and not its correctness.
% AFIRMACAO: A33
In S951, with GPT4omini on 67 samples, the VIR metric (lower is better) was 100.0\% in the baseline, 74.6\% with security requirements and 67.2\% with retrieval of vulnerable cases, with no explanation for the baseline of 100\% and with functionality measured on another set (S951-C1; S951-C2; S951).
This result on 67 samples does not demonstrate security in real use.

\textbf{Prompt optimization and hybrid systems.}
% AFIRMACAO: A34
In S650, with GPT-3.5, automatic prompt optimization raised Pass@1 from 60.5\% to 73.0\% from Java to Python on CodeNet, with one test case per task, and the ablation of the iteration was made only on a generation task (S650-C2).
S50 includes components outside the scope, and S1220 and S1836 are descriptive.

\textbf{Gaps.} The evidence is insufficient on the following points:
\begin{itemize}
  \item no extracted comparison isolates the reasoning from the example; S193 isolates the explanation, but with one model and one pair of languages;
  \item prompt optimization for translation has 1 study with 1 comparison (S650);
  \item security has 1 study (S951), with one model in the extracted comparisons;
  \item migration of versions and domains (Java 8 to 11, Python to Spark, SQL, VHDL, Triton) only appears in studies without a counterfactual, cumulative or with an internal set (S1220, S1836, S1803, S1132);
  \item no translation comparison measures cost or number of calls.
\end{itemize}

\textbf{Studies extracted after the partial synthesis.}
% AFIRMACAO: A115
In S1713, with Claude-2 in the translation of 569 C++ programs to Rust, \textit{few-shot} raised from 240 to 539 the
programs that compile. The counterexamples added nothing beyond \textit{few-shot} (539), and the package includes
repair by manual rules (S1713-C1 to S1713-C3).
% AFIRMACAO: A116
In S1807, with Llama 3.1 70B translating C programs from CodeNet, curated \textit{zero-shot} changed the rate of
translations that do not compile from 14.91\% to 14.79\%, and CoT raised it to 21.36\% (S1807-C1; S1807-C2).

\subsubsection{Refactoring}

The partial synthesis of refactoring covers 19 refactoring studies, and the one included later is at the end of this part. With overlap between the types, 12 of the 19 studies evaluate at least one package, 9 isolate some ingredient, 1 has a cumulative design (S2155) and 4 have no counterfactual (S739, S1368, S1698 and S1706). In S958, the baseline is the human code from LeetCode, and not another prompting condition. Fifteen of the 19 studies use their own sets or specific projects, and contamination is recorded as not reported in 13 of the 19. The groupings are:
\begin{itemize}
  \item examples and retrieval of examples (6 studies);
  \item specification of the transformation and domain knowledge (6);
  \item explicit reasoning (2);
  \item feedback from tests or compilation, self-repair and iteration (5);
  \item studies without a counterfactual (4).
\end{itemize}

\textbf{Examples.}
% AFIRMACAO: A35
With examples and the same model in both conditions, the test pass rate went from 28.36\% to 34.51\% in S4 (StarCoder2-15B, \textit{one-shot}), Pass@1 went from 87.44\% to 91.11\% in S202 (gpt-3.5-turbo-0301) and correctness went from 87.33\% to 90.0\% in S1822 with one retrieved example (DeepSeek-V3), in the row of structure distortions D1 (S4-C2; S202-C1; S1822-C1).
% AFIRMACAO: A36
There are also side effects or null results: in S202, compilability fell from 99.99\% to 99.93\% and the cyclomatic complexity of \textit{one-shot} got worse; in S1685, the compilations went from 31 to 32 in 37 samples, without a statistical test; and in S958 the execution time went from 0.00606 to 0.00617 ms, in comparison with the human code (S202; S1685-C2; S958-C4).

\textbf{Specification of the transformation.} The six studies name the type of refactoring or inject schemas, changelogs, taxonomies, rules or definitions; S466, S511, S799 and S2122 isolate ingredients.
% AFIRMACAO: A37
In S511, with GPT-4 on 180 cases, in cumulative stages, specifying the type of refactoring raised recall from 15.6\% to 52.2\%, adding subcategory and example raised it to 66.7\%, and limiting the search space took it from 66.7\% to 86.7\%, a value that the conclusion of the study records as 86.1\% (S511-C1; S511-C2; S511-C3; S511).
% AFIRMACAO: A38
In S2122, the package raised Pass@1 from 4\% to 76\% with Mistral-Large, and with LLaMA-4-Maverick removing the API schema or the changelogs gave 56\% and 60\%, against 76\% for the package, but CrystalBLEU fell from 0.78 to 0.66 while Pass@1 rose (S2122-C1; S2122-C2; S2122-C3; S2122).
% AFIRMACAO: A39
In S466, the reduction of smells went from 55\% to 89\%, and in S799 the true positives went from 10 to 12 in 25 snippets, with one invocation; the gain of the rule DSL in S466 (from 67\% to 89\%) was reported only in prose, with the model not reported (S466-C1; S466-C2; S799-C1).
% AFIRMACAO: A40
In the opposite direction, in S362 a rule-based baseline (RIdiom) had higher precision than the approach of the study in eight of the nine idiomatic Python patterns, by 0.8 to 6.2 percentage points (Figure 6c) (S362-C1; S362).
One possible reading is that the largest gains appear when the baseline is very low and the task depends on version knowledge (S511-C2; S2122-C1).

\textbf{Explicit reasoning.}
% AFIRMACAO: A41
Explicit reasoning raised the pass rate from 28.36\% to 32.22\% in S4, evaluated as a package, and in S958 the energy went from 1.7122 to 1.7112 mWh in comparison with the human code, which the study itself describes as almost equal consumption (S4-C1; S958-C1; S958).
Measuring energy on a benchmark does not demonstrate a reduction in operating cost.

\textbf{Feedback and iteration.}
% AFIRMACAO: A42
Among the favourable results, the TPP proportion went from 0.503 to 0.638 in S1054 (CodeQwen-1.5), with a comparator that also differs in the example and with a narrative that contradicts the table, syntactic validity went from 86.1\% to 92.4\% in S2183 and the compilations went from 23 to 32 in S1685, against GPT with an undescribed prompt (S1054-C1; S2183-C1; S1685-C1).
% AFIRMACAO: A43
There are also reversals: in S1685, the condition without the RefactoringMiner feedback had more compilations (34 against 32 for the package), and in S2164, with GPT-4, the second round raised acceptance in Python on LeetCode from 72\% to 73.53\%, whereas, in the Easy Before subset, asking for optimization increased the time from 86.24 ms to 90.53 ms (S1685-C3; S2164-C2; S2164-C1).
% AFIRMACAO: A44
In S739, after 10 iterations with GPT-4o without human intervention, the vulnerability totals observed over the iterations were 124 in the efficiency strategy, 158 in the functionality one, 67 in the generic improvement one and 38 in the security one; the study compares the strategies with one another, without a control without iterative improvement (S739-C1; S739-C2; S739-C3).
In S1685, the model is not reported and reproducibility is not met; in S739, the comparisons are between strategies, and not against a condition without the technique.

\textbf{Gaps.} The evidence is insufficient on the following points:
\begin{itemize}
  \item isolated reasoning: the 2 studies evaluate it as a package or against human code;
  \item tests, architecture, Terraform, SQL and decompiled code, with 1 study each (S466, S1775, S2122, S2183 and S1822);
  \item adoption by developers, addressed only by S1368, which is descriptive.
\end{itemize}

\textbf{Study extracted after the partial synthesis.}
% AFIRMACAO: A117
S1883 reports mean speedups of 2, 2.95 and 3.7 times on 2, 4 and 8 nodes when converting, with GPT-4o, the 100
most time-consuming serial codes to distributed execution, in comparison with serial execution. The result is that of a
system that combines prompting and execution environment, and there is no comparator that isolates the prompting (S1883-C1).
